\documentclass[11pt]{article}

\usepackage[margin=1in]{geometry}
\usepackage{amsmath, amssymb}
\usepackage{setspace}
\usepackage{hyperref}

\setstretch{1.1}

\title{\textbf{STA380 Project Proposal}\\
Monte Carlo Study of the Bias--Variance Tradeoff in Regression Models}

\author{
Jizheng Huang \\
\texttt{jizheng.huang@mail.utoronto.ca}
\and
Victor Jiang \\
\texttt{vic.jiang@mail.utoronto.ca}
\and
Tianchen Xu \\
\texttt{tianchen.xu@mail.utoronto.ca}
\and
Kai Rui Zhu \\
\texttt{kairui.zhu@mail.utoronto.ca}
}

\date{University of Toronto, Mississauga}

\begin{document}
\maketitle

% ---------------------------
\section{Project Topic}
Monte Carlo Study of the Bias--Variance Tradeoff in Regression Models.


% ---------------------------
\section{Simulation vs.\ Dataset}
This project will use pure simulation rather than a real-world dataset.

Data will be generated according to the model
\[
y = f(x) + \varepsilon,
\]
where $f(x)$ is a known nonlinear regression function and $\varepsilon$ is a random noise term. Knowing the true data-generating process allows us to estimate bias and variance using repeated Monte Carlo experiments.

% ---------------------------
\section{Project Details}
Monte Carlo simulation will be used to repeatedly generate independent training and test datasets. Regression models of varying complexity will be fitted to the training data and evaluated on test data.

The following regression methods will be considered:
\begin{itemize}
    \item Polynomial regression, with model complexity controlled by the polynomial degree
    \item $k$-nearest neighbor (kNN) regression, with complexity controlled by the number of neighbors $k$ 
\end{itemize}

For each model and complexity level, we will compute:
\begin{itemize}
    \item Training mean squared error (MSE)
    \item Test mean squared error (MSE)
    \item Monte Carlo estimates of bias$^2$ and variance
\end{itemize}

Bias and variance will be estimated by repeatedly fitting models to independently generated datasets and examining the systematic deviation and variability of predictions relative to the true regression function. Results will be summarized across different sample sizes, noise levels, and model complexities to illustrate overfitting and underfitting behavior.

% ---------------------------
\section{User Inputs (Shiny Components)}
The R Shiny application will include at least the following interactive inputs:
\begin{enumerate}
    \item Sample size for the simulated dataset
    \item Noise level of the error term $\varepsilon$
    \item Random seed for reproducibility
    \item Choice of regression model (polynomial regression or kNN regression)
    \item Model complexity parameter (polynomial degree or number of neighbors $k$)
    \item Option to download simulated data or results as a \texttt{.csv} file
\end{enumerate}

Additional inputs may be included to enhance interactivity, such as plot customization options or the ability to compare multiple model complexities simultaneously.

% ---------------------------
\section{Generative AI Usage Statement}
We used generative AI tools to assist in writing the LaTeX code for the proposal. All statistical methods, project design decisions, and final content were reviewed and validated by team members. We did not use any generative AI tools to generate results, simulations, or analyses.

% ---------------------------

\bibliographystyle{plain}
\nocite{*}
\bibliography{references}

\end{document}
